\printedpage{4}
% head: A. Grothendieck : Dérivateurs, Ch. XI, Édition des Universités de Montpellier II et Paris VII
% foot: Version 28 mai 2002

\textbf{Proposition 1.} \textit{This notion is stable under base change (trivial), under composition (formal from a)). For $f$ to be a \textup{Hot}-fibration, it is necessary and sufficient that $f$ \textup{\texttt{[rather}} $f^{o}$\textup{\texttt{]}} be one.}

\textbf{Proposition 2.} \textit{Let $f : X \longrightarrow S$ be a \textup{Hot}-fibration. Then $s \longmapsto \mathrm{hot}(X_s)$ is a local system on $S$, with values in \textup{Hot}. Thus the $X_s$ for $s$ in the same connected component of $S$ are \textup{Hot}-isomorphic to one another (but not canonically, if this component is not \textup{1}-connected).}

% (a vertical change bar is printed in the left margin beside Théorème 2)
\textbf{Theorem 2.} {\itshape Let
\begin{tikzcd}[ampersand replacement=\&, baseline=(\tikzcdmatrixname-1-1.base), column sep=small, row sep=small]
X' \arrow[rr, "f"] \arrow[dr, "p'"'] \& \& X \arrow[dl, "p"] \\
\& S \&
\end{tikzcd}
be a commutative diagram in \textup{Cat}, with $p$, $p'$ \textup{Hot}-fibrations, and let $S_0 \subseteq S$ be a subset of $S$ such that $S_0 \longrightarrow \pi_0(S)$ is surjective. For $f$ to be a \textup{Hot}-equivalence, it is necessary and sufficient that the $f_s : X_s \longrightarrow Y_s$ $(s \in S_0)$ be so.} % [sic] f_s : X_s -> Y_s -> f_s : X'_s -> X_s, the diagram has X' and X (cf. f_s : X'_s -> X_s p. 6)

a) Sufficiency is proved as before, with the Artin-Mazur criterion.

\msp{page 132}

Let us note that the hypothesis on $S_0$ ensures that if $F$, $F'$ are two locally constant sheaves of sets on $S$, and $u : F \longrightarrow F'$ a homomorphism such that $u_s : F_s \xrightarrow{\ \sim\ } F'_s$ for $s \in S$, then $u$ is \texttt{[an]} isomorphism. % [sic] s \in S -> s \in S_0, otherwise the hypothesis on S_0 plays no role

1$^{o}$) Let $F$ be a locally constant sheaf of sets on $X$, $F'$ its inverse image on $X'$, we must prove
\[
H^0(X,F) \ \xrightarrow{\ \sim\ }\ H^0(X',F').
\]
Now
\[
H^0(X,F) \simeq H^0(Y, \underbrace{f_*(F)}_{\substack{\text{locally}\\ \text{constant on } S}}) \ \longrightarrow\ H^0(Y, f'_*(F')) \simeq H^0(X',F'),
\] % [sic] H^0(Y, f_*(F)), H^0(Y, f'_*(F')) -> H^0(S, p_*(F)), H^0(S, p'_*(F')), the brace says "localement constant sur S" and (*) uses p_*, p'_*
we must prove that
\[
p_*(F) \ \xrightarrow{\ \sim\ } p'_*(F') \tag*{$(*)$}
\]
is \texttt{[an]} isomorphism. Now by the hypothesis on $p$, $p'$, the two sheaves are locally constant and are computed fiber by fiber, from which it follows that $(*)$ is an isomorphism for the fibers at $s \in S_0$, hence an isomorphism everywhere. OK.

2$^{o}$) It already follows from this, if $F$ is a sheaf of \textit{groups,} that the functor
\[
(F\text{-torsors on } X) \ \longrightarrow\ (F\text{-torsors on } X')
\]
is fully faithful, hence $H^1(X,F) \longrightarrow H^1(X',F')$ \texttt{[is]} injective; it remains to prove that this functor is essentially surjective. Because of full faithfulness, this is \textit{local on $S$} (\textbf{NB} We see directly that the hypothesis made on the $s \in S_0$ implies the same hypothesis for every $s \in S$, using the fact that $p$, $p'$ are fibrations) (\footnote{The hypotheses made are indeed local on $S$ (at least if $S_0 = S$).}). Thus we are reduced to the case where $S = S/s_0$ has a final object $s_0$.

\printedpage{5}
% head: A. Grothendieck : Dérivateurs, Ch. XI, Édition des Universités de Montpellier II et Paris VII
% foot: Version 28 mai 2002

\msp{page 133}

We then have
\[
\begin{aligned}
H^1(X,F) &\ \simeq\ H^0(S, R^1p_*(F)) \\
H^1(X',F') &\ \simeq\ H^0(S, R^1p'_*(F')),
\end{aligned}
\]
and we see again that
\[
R^1p_*(F) \ \longrightarrow\ R^1p'_*(F')
\]
is \texttt{[an]} isomorphism (for it is computed fiber by fiber, and we use the hypothesis on the $f_s$), hence induces an isomorphism
\[
H^1(X,F) \ \longrightarrow\ H^1(X',F').
\]

3$^{o}$) If $F$ is even abelian, we must prove
\[
H^i(X,F) \ \xrightarrow{\ \sim\ }\ H^i(X',F')
\]
for every $i$, and this again follows from the Leray spectral sequence for $(p,F)$ and for $(p',F')$.

b) Conversely, suppose that $f$ is a Hot-equivalence, and let us prove that the $f_s : X'_s \longrightarrow X_s$ are Hot-equivalences. I am going to use the

% (a vertical change bar is printed in the left margin beside the Lemme)
\textbf{Lemma.} \textit{Let $S \in \Ob\Cat$, $S$ \textup{0}-connected. Then there exists an $\overline{S} \in \Ob\Cat$, $\overline{S}$ \textup{Hot}-aspheric, and $\overline{S} \longrightarrow S$ a \textup{Hot}-fibration.}

Let us see how this implies theorem 2 (converse part). We may suppose $S$ 0-connected, and we make the base change by $\overline{S} \longrightarrow S$

\msp{page 134}

as in the lemma:
\[
\begin{tikzcd}[column sep=small, row sep=small]
X' \arrow[rr, "f"] \arrow[dr, "p'"'] & & X \arrow[dl, "p"] & & & & \\
& S & & & \overline{X'} \arrow[rr, "\overline{f}"] \arrow[dr, "\overline{p'}"'] & & \overline{X} \arrow[dl, "\overline{p}"] \\
& & & & & \overline{S} \arrow[ullll, "\varphi"] &
\end{tikzcd}
\quad .
\]

Then $\overline{p'}$ and $\overline{p}$ are Hot-fibrations, and since $\varphi$ is a Hot-fibration, base change by $\varphi$ transforms the Hot-equivalence $f$ into a Hot-equivalence $\overline{f}$. Consider then $\overline{s} \in \Ob\overline{S}$, and the diagram
\[
\begin{tikzcd}
\overline{X'} \arrow[r, "\overline{f}"] & \overline{X} \\
\overline{X'}_{\overline{s}} \arrow[u, hook, "i'"] \arrow[r, "\overline{f}_{\overline{s}}"] & \overline{X}_{\overline{s}} \arrow[u, hook, "i"']
\end{tikzcd}
\qquad .
\]

\printedpage{6}
% head: A. Grothendieck : Dérivateurs, Ch. XI, Édition des Universités de Montpellier II et Paris VII
% foot: Version 28 mai 2002

Since $\{\overline{s}\} \longrightarrow \overline{S}$ is a Hot-equivalence, the arrow $i$ deduced from it by the base change $\overline{p} : \overline{X} \longrightarrow \overline{S}$ is one as well, since $\overline{p}$ is a Hot-fibration, and likewise for $i'$. Thus in the preceding commutative square, $\overline{f}$, $i$ and $i'$ are Hot-equivalences, hence so is $\overline{f}_{\overline{s}}$. But if $s = \varphi(\overline{s})$, $\overline{f}_{\overline{s}}$ is identified with $f_s : X'_s \longrightarrow X_s$. Besides, it is immediate that $\varphi$, being a Hot-fibration

\msp{page 135}

with nonempty source, and connected target, is surjective. (\textbf{NB} all the fibers of one and the same connected component of $S$ are Hot-equivalent, hence if one is nonempty, they all are.) This completes the proof, modulo the lemma.

It remains to prove it. Since $S = \emptyset$ \texttt{[rather} $S \neq \emptyset$\texttt{]}, I can choose an ``origin''
\[
s_0 \in S,
\]
and look at the category
\[
\Ch_{\infty}(S; s_0, -)
\]
of infinite paths in $S$, with origin $s_0$, fiber of $\Ch_{\infty}(S)$ under the source functor
\[
\sigma : \Ch_{\infty}(S) \ \longrightarrow\ S.
\]
We know that the fibers of $\sigma$ are Hot-aspheric (VII), and the category constructed here is none other than (with the notations of VII, p.~152) than the category $\overline{S'}$ which enters into the canonical factorization of
\[
f : \{s_0\} = S' \ \longrightarrow\ S
\]
into
\[
\{s_0\} = S' \ \underbrace{\longrightarrow}_{\text{Hot-equivalence}}\ \overline{S'} \ \xrightarrow{\ \overline{f}\ }\ S,
\]
where here $\overline{f}$ is interpreted simply as the target functor (i.e.\ the functor induced on

\msp{page 136}

$\Ch_{\infty}(S; s_0, -)$ by the target functor
\[
\beta : \Ch_{\infty}(S) \ \longrightarrow\ S \ \ ).
\]
We proved in loc.\ cit.\ that $\overline{f}$ is a Hot-fibration (``proper'' and ``smooth''), but in a context which seemed to appeal tacitly to the theory of the derivator HOT. So to put the ``topological'' results of loc.\ cit.\ on a solid basis, we must develop the relevant notions of proper functor, smooth functor, without reference to a HOT theory. This is what I had in fact already done in 1983, in the days of Pursuing Stacks, but without a proper write-up on this subject. So I must now return to these notions, in a somewhat systematic way. Besides, I would have trouble constructing a single Hot-fibrant functor (except for the trivial examples, such as isomorphisms, or the products $Y \times Z \xrightarrow{\mathrm{pr}_1} Y$), if I did not have at my disposal a theory of smooth functors and proper functors (\footnote{I am exaggerating a little all the same -- after all it must be fairly obvious that the Thomason-Kan functors are Hot-fibrations (but there are still a few too few of them!).}). % [sic] dû mal -> du mal, "avoir du mal à" (partitive article, no accent)